% \subsubsection{CSA Two-Stage Calibration}
% \begin{itemize}
%     \item Training vectors are split into two disjoint halves $S_1$ and $S_2$ (50/50 random split).
%     \item $S_1$ (shape discovery): used to learn the geometry of the envelope.
%     \item $S_2$ (size scaling): used to calibrate the scaling factor $\hat{t}$ that determines the final size of the envelope.
%     \item The two stages must use disjoint data to preserve the exchangeability argument underlying the coverage guarantee.
% \end{itemize}

% \subsubsection{\texorpdfstring{$\tau$ Computation}{tau Computation}}
% \begin{itemize}
%     \item Define $\tau(\mathbf{s}) = \inf\{t > 0 : \mathbf{s} \in t \cdot E\}$, the minimum scaling factor needed to include point $\mathbf{s}$ in the scaled envelope $E$.
%     \item The global scaling factor is then determined by the $(1-\alpha)$ quantile of the $\tau$ values over $S_2$:
%     \[
%     \hat{t} = \tau_{(k)} \quad \text{where} \quad k = \lceil(|S_2|+1)(1-\alpha)\rceil
%     \]
%     where $\tau_{(k)}$ denotes the $k$-th smallest value of $\tau$ in $S_2$.
%     \item Key contribution: for each of the three envelope geometries, $\tau(\mathbf{s})$ has a closed-form expression computable in a single pass — no iterative search required.
%     \item Computational benefit: the entire calibration step reduces to one forward pass over $S_2$ followed by a sort.
%     \item Contrast with the original CSA paper which uses iterative binary search.
% \end{itemize}


\subsection{Envelope Construction}
\label{sec3.4: envelope construction}
\noindent Here we will present the three envelope geometries that we have evaluated in this report. Each is defined by a 
calibration procedure on $S_1$ that fixes the shape of $E$, and a test time rule for evaluating $\tau(\mathbf{s}) = 
\inf\{t>0 : \mathbf{s} \in t\cdot E\}$ in closed form. This replaces the iterative search that was used in the original CSA 
method (Section~\ref{sec2.4.2: csa framework}) with a single forward pass over $S_2$ followed by a sort.

\subsubsection{Collapsed 1D Envelope Method}
\label{sec3.4.1: collapsed 1d}
\noindent The simplest envelope construction method just collapses the $K$-dimensional (or $K_h$-dimensional) NC vector to a 
scalar value by taking the mean over the over observed dimensions,
\begin{equation*}
\bar{s} = \frac{1}{\|\mathbf{m}\|_1} \sum_{k : m_k = 1} v_k,
\end{equation*}
\noindent and calibrates a single quantile on $S_1$:

\begin{equation*}
\tilde{q} = \operatorname{Quantile}_{1-\alpha}\big(\{\bar{s}_i\}_{i \in S_1}\big), \qquad 
\tau(\mathbf{s}) = \frac{\bar{s}}{\tilde{q}}, \qquad E = \{\mathbf{s} : \bar{s} \le \tilde{q}\}.
\end{equation*}

\noindent This envelope is star-shaped (Section~\ref{sec3.2: non-convex envelopes}), it is unaffected by the dimensionality 
$K$, but at the same time is also discards all the geometric information that the $K$ dimensions provide us with. Any two 
vectors with the same mean are treated identically regardless of the individual coordinates.


%%%%%%%%%%%%%%%%%%%%%%%%%%%%%%%%%%%

\subsubsection{Radial Envelope Method}
\label{sec3.4.2: radial}
\noindent The Radial envelope method retains the directional structure of the scores by calibrating a boundary radius as a 
function of direction on the unit sphere which is restricted only to the positive orthant (since all NC values are 
non-negative). In this method, we draw $M$ directional vectors $\{\mathbf{u}_m\}_{m=1}^M$ uniformly over the orthant by 
sampling $\mathbf{v}_m \sim |\mathcal{N}(0, I_K)|$ coordinate-wise and normalizing: $\mathbf{u}_m = \mathbf{v}_m / \|
\mathbf{v}_m\|$.

\vspace{0.15cm}
\noindent For each $S_1$ point $\mathbf{s}_i$, over the observed dimensions only, we compute its magnitude and unit direction, 
$r_i = \|\mathbf{s}_{i,\mathbf{m}_i}\|$, $\mathbf{d}_i = \mathbf{s}_{i,\mathbf{m}_i}/r_i$. For each sample direction 
$\mathbf{u}_m$, we collect $n_m$ calibration points within an angular sector of half-width $\delta = 30^\circ$ (where 
$\mathbf{d}_i \cdot \mathbf{u}_m \ge \cos\delta$). The directional radius $\tilde{q}_m$ is then calibrated as:
\begin{equation*}
\tilde{q}_m = \begin{cases}
\operatorname{Quantile}_{1-\alpha}\big(\{r_i : \mathbf{d}_i \cdot \mathbf{u}_m \ge \cos\delta\}\big) & \text{if } 
n_m \ge n_{\min}, \\
w_m \cdot \operatorname{Quantile}_{1-\alpha}\big(\{r_i : \mathbf{d}_i \cdot \mathbf{u}_m \ge \cos\delta\}\big) + 
(1-w_m) \cdot \tilde{q}_{\text{marg}} & \text{if } 0 < n_m < n_{\min}, \\
\tilde{q}_{\text{marg}} & \text{if } n_m = 0,
\end{cases}
\end{equation*}

\noindent where $n_m$ is the number of $S_1$ points within the band of direction $m$, $n_{\min}$ a minimum-support threshold, $w_m = 
n_m/n_{\min}$, and $\tilde{q}_{\text{marg}} = \operatorname{Quantile}_{1-\alpha}(\{r_i\}_{i \in S_1})$ the marginal fallback 
radius. This blend enables us to avoids a hard cutoff between sufficient and insufficient local support (number of points in 
that region).

\vspace{0.15cm}
\noindent To ensure that we get a smooth boundary, we compute the radius at an arbitrary direction $\mathbf{d}$ as a 
directional weighted average of all $M$ calibrated radii:
\begin{equation}
\label{eq: radial boundary}
r(\mathbf{d}) = \sum_{m=1}^M w_m(\mathbf{d}) \, \tilde{q}_m, \qquad 
w_m(\mathbf{d}) = \frac{\exp(\kappa \, \mathbf{d}\cdot\mathbf{u}_m)}{\sum_{m'=1}^M \exp(\kappa \, \mathbf{d}\cdot \mathbf{u}_{m'})},
\end{equation}
where $\kappa > 0$ controls the angular smoothness of the interpolation across directions (larger $\kappa$ approaches 
nearest-direction assignment). This makes the envelope boundary a continuous function of direction rather than a piecewise 
constant one. The scaling function follows directly from this:
\begin{equation*}
\tau(\mathbf{s}) = \frac{\|\mathbf{s}_{\mathbf{m}}\|}{r(\mathbf{d})}, \qquad \mathbf{d} = 
\frac{\mathbf{s}_{\mathbf{m}}}{\|\mathbf{s}_{\mathbf{m}}\|}, \qquad E = \{\mathbf{s} : \|\mathbf{s}_{\mathbf{m}}\| \le 
r(\mathbf{d})\}.
\end{equation*}
As established in Section~\ref{sec3.2: non-convex envelopes}, this construction is star-shaped by design. This is because $r$ 
depends only on direction and not on magnitude, $\mathbf{s} \in t\cdot E$ for some $t$ implies $\lambda\mathbf{s} \in t\cdot E$ 
for all $\lambda \in [0,1]$.


%%%%%%%%%%%%%%%%%%%%%%%%%%%%%%%%%%%%%%%%%%%%%%%%%%%%%%%%%%%%%%%%%%%%%%%%%%%%%%%%%%%%%%%%%%%%%%%%%%%%%

\subsubsection{Strip Envelope Method}
\label{sec3.4.3: strip}
The Strip envelope defines conditional bounds for each target dimension $i$ based on the bin occupied by another dimension 
$j$. We partition the range of dimension $j$ into $NB$ equal-width bins. For each bin $n$ containing $n_{j,n}$ calibration 
points, the upper limit along dimension $i \neq j$ is calibrated as:
\begin{equation*}
\text{limits}[j,i,n] = \begin{cases}
\operatorname{Quantile}_{1-\alpha}\big(\{s_{\cdot,i} : s_{\cdot,j} \in \text{bin } n\}\big) & \text{if } n_{j,n} \ge n_{\min}, \\
w_{j,n} \cdot \operatorname{Quantile}_{1-\alpha}\big(\{s_{\cdot,i} : s_{\cdot,j} \in \text{bin } n\}\big) + (1-w_{j,n}) \cdot 
\tilde q_i & \text{if } 0 < n_{j,n} < n_{\min}, \\
\tilde{q}_i & \text{if } n_{j,n} = 0,
\end{cases}
\end{equation*}
\noindent where $\tilde{q}_i$ is the marginal $(1-\alpha)$ quantile for dimension $i$ over $S_1$ , and $w_{j,n} = n_{j,n} / n_{\min}$ 
shrinks underpopulated bins toward it (it mirrors is the blending rule of the radial envelope).

\vspace{0.15cm}
\noindent A windowed monotonicity pass then enforces $\text{limits}[j,i,n] \ge \text{limits}[j,i,n']$ for $n' \in (n, n+w]$ 
within a fixed window $w$. This is done to ensure that a bin's limit cannot be undercut by a bin immediately ahead of it in 
the same conditioning dimension. \todo{explain why it shouldnt happen}

\vspace{0.15cm}
\noindent The scaling function $\tau(\mathbf{s})$ aggregates the worst-case ratio over all conditioning pairs $(j,i)$:
\begin{equation*}
\tau(\mathbf{s}) = \max_{j \neq i} \frac{s_i}{\text{limits}[j,i,\text{bin}_j(\mathbf{s})]}, \qquad E = \{\mathbf{s} : 
\forall j\neq i,\; s_i \le \text{limits}[j,i,\text{bin}_j(\mathbf{s})]\}.
\end{equation*}

\noindent Missing dimensions ($m_k = 0$) are handled pairwise: 
\begin{itemize}
    \item \textbf{Conditioning dimension $j$ missing:} The bin cannot be determined so the ratio falls back to dimension $i$'s 
           marginal quantile $\tilde{q}_i$ instead of the conditional limit.
    \item \textbf{Target dimension $i$ missing:} Its value $s_i$ is replaced by its marginal qunatile $\tilde{q}_i$ before
           computing the ratio.
    \item \textbf{Both missing:} The pair contributes no constraint ($\text{ratio} = 0$), since absence of evidence should not 
          lead to a false rejection.
\end{itemize}

\subsubsection{Class-Specific \texorpdfstring{$\hat{t}$}{t̂} Calculation}



%%%%%%%%%%%%%%%%%%%%%%%%%%%%%%%%%%%%%%%%%%%%%%%%%%%%%%%%%%%%%%%%%%%%%%%%%%%%%%%%%%%%%%%%%%%%%%%%%

\vspace{4cm}



\subsubsection{Class-Conditional Calibration}
\noindent For any of the three constructions above, a single scaling factor computed over all of $S_2$ yields only 
marginal coverage (Section~\ref{sec2.3.3: coverage types}). We instead compute a per-class threshold and take the 
maximum:
\begin{equation}
\label{eq: class conditional that}
\hat{t}_h = \tau_{\lceil (|S_2^h|+1)(1-\alpha)\rceil}, \qquad \hat{t} = \max_{h \in \mathcal{H}} \hat{t}_h,
\end{equation}
where $S_2^h = \{\mathbf{s} \in S_2 : H = h\}$ and $\tau_{(k)}$ denotes the $k$-th order statistic of $\{\tau(\mathbf{ 
s})\}_{\mathbf{s}\in S_2^h}$. Taking the maximum over classes guarantees $\ge 1-\alpha$ coverage for every class 
individually, at the cost of set size for classes whose $\hat t_h$ is below the maximum. The added computational cost 
over the marginal $\hat t$ is negligible: a per-class sort of $\tau$-values already computed in one pass.

\subsubsection{Resolving the Empty Prediction Set}
\noindent Section~\ref{sec3.1: problem formalisation} noted that $C(\mathbf{s},\mathbf{m})$ may be empty when no 
candidate label's transformed score falls inside $\hat t \cdot E$. The two tasks resolve this differently. For 
gram-type classification, an empty set is resolved conservatively by returning both labels, $C(\mathbf{s}) = 
\mathcal{H}$, rather than an arbitrary tie-break — since $N=2$, this is the maximally uninformative but always-valid 
choice. For host-level classification, where $N$ can be large and returning all of $\hat{\mathcal{H}}(\mathbf{s})$ 
would be uninformative, an empty set is instead resolved by forcing a singleton at the host with minimum $\tau$:
\begin{equation*}
C(\mathbf{s}) = \Big\{ \operatorname*{argmin}_{h \in \hat{\mathcal{H}}(\mathbf{s})} \tau_h(\mathbf{s}_h) \Big\}.
\end{equation*}
Both resolutions are applied only when the primary rule of Equation~\eqref{eq: prediction set definition} yields 
$\emptyset$; the coverage guarantee of Theorem~\ref{thm: marginal coverage} is unaffected, since it is established on 
the primary construction over $\mathcal{H}$ and empty outcomes are, by definition, cases where that guarantee already 
requires \emph{some} label to be returned to satisfy the $1-\alpha$ bound in expectation. We note, however, that 
neither task's reported \texttt{empty\_rate} can be nonzero as a result — this reflects the resolution rule, not an 
absence of out-of-envelope test points, and Section~\ref{sec5: experiments} reports the pre-resolution 
\texttt{forced\_rate} separately to make this visible.